\ifdefined\gkver\else\def\gkver{student}\fi
\documentclass[\gkver]{gaokaozhenti}
\begin{document}

\chapter{2015年江苏理综}
%% paperType: 理综物理部分

\begin{enumerate}
\item
%% number: 1
%% paperName: 2015年江苏理综第 1 题 3 分
%% typeId: 1
%% type: 单选题
%% chapter: 交流电
%% point: 变压器
%% method: 无
%% score: 3
%% degree: 850
%% duplicateId: 0
%% body:
一电器中的变压器可视为理想变压器，它将 $220\UV$ 交变电流改变为 $110\UV$。已知变压器原线圈匝数为 800，则副线圈匝数为\xzanswer{B}

\fourchoices[answer=B]
{200}
{400}
{1600}
{3200}

%% answer: B
%% memo:
\memoanswer{
	本题原卷及材料中未提供详解，待补充。
}

\item
%% number: 2
%% paperName: 2015年江苏理综第 2 题 3 分
%% typeId: 1
%% type: 单选题
%% chapter: 电场
%% point: 静电
%% method: 无
%% score: 3
%% degree: 650
%% duplicateId: 0
%% body:
静电现象在自然界中普遍存在，我国早在西汉末年已有对静电现象的记载，《春秋纬$\cdot$考异邮》中有“玳瑁吸菸”之说，但下列不属于静电现象的是\xzanswer{C}

\fourchoices[answer=C]
{梳过头发的塑料梳子吸起纸屑}
{带电小球移至不带电金属球附近，两者相互吸引}
{小线圈接近通电线圈过程中，小线圈中产生电流}
{从干燥的地毯上走过，手碰到金属把手时有被电击的感觉}

%% answer: C
%% memo:
\memoanswer{
	本题原卷及材料中未提供详解，待补充。
}

\item
%% number: 3
%% paperName: 2015年江苏理综第 3 题 3 分
%% typeId: 1
%% type: 单选题
%% chapter: 曲线运动
%% point: 天体运动
%% method: 无
%% score: 3
%% degree: 650
%% duplicateId: 0
%% body:
过去几千年来，人类对行星的认识与研究仅限于太阳系内，行星“51 peg b”的发现拉开了研究太阳系外行星的序幕。“51 peg b”绕其中心恒星做匀速圆周运动，周期约为 4 天，轨道半径约为地球绕太阳运动半径的 1/20。该中心恒星与太阳的质量比约为\xzanswer{B}

\fourchoices[answer=B]
{$\frac{1}{10}$}
{$1$}
{$5$}
{$10$}

%% answer: B
%% memo:
\memoanswer{
	本题原卷及材料中未提供详解，待补充。
}

\item
%% number: 4
%% paperName: 2015年江苏理综第 4 题 3 分
%% typeId: 1
%% type: 单选题
%% chapter: 磁场
%% point: 左手定则
%% method: 等效替代
%% score: 3
%% degree: 450
%% duplicateId: 0
%% body:
如图所示，用天平测量匀强磁场的磁感应强度。下列各选项所示的载流线圈匝数相同，边长 MN 相等，将它们分别挂在天平的右臂下方。线圈中通有大小相同的电流，天平处于平衡状态。若磁场发生微小变化，天平最容易失去平衡的是\xzanswer{A}

\onepicture[h=3.6cm]{figs/04.png}

\fourchoices[answer=A, ispicture=true, h=2.2cm]
{figs/04a.png}
{figs/04b.png}
{figs/04c.png}
{figs/04d.png}

%% answer: A
%% memo:
\memoanswer{
	本题原卷及材料中未提供详解，待补充。
}

\item
%% number: 5
%% paperName: 2015年江苏理综第 5 题 3 分
%% typeId: 1
%% type: 单选题
%% chapter: 直线运动
%% point: 多段匀变速直线运动
%% method: 无
%% score: 3
%% degree: 450
%% duplicateId: 0
%% body:
如图所示，某“闯关游戏”的笔直通道上每隔 $8\Um$ 设有一个关卡，各关卡同步放行和关闭，放行和关闭的时间分别为 $5\Us$ 和 $2\Us$。关卡刚放行时，一同学立即在关卡 1 处以加速度 $2\Umsq$ 由静止加速到 $2\Ums$，然后匀速向前，则最先挡住他前进的关卡是\xzanswer{C}

\onepicture[h=2.6cm]{figs/05.png}

\fourchoices[answer=C]
{关卡 2}
{关卡 3}
{关卡 4}
{关卡 5}

%% answer: C
%% memo:
\memoanswer{
	本题原卷及材料中未提供详解，待补充。
}

\item
%% number: 6
%% paperName: 2015年江苏理综第 6 题 4 分
%% typeId: 2
%% type: 多选题
%% chapter: 运动和力
%% point: 超重失重
%% method: 无
%% score: 4
%% degree: 650
%% duplicateId: 0
%% body:
一人乘电梯上楼，在竖直上升过程中加速度 $a$ 随时间 $t$ 变化的图线如图所示，以竖直向上为 $a$ 的正方向，则人对地板的压力\xzanswer{AD}

\onepicture[h=3.4cm]{figs/06.png}

\fourchoices[answer=AD]
{$t=2\Us$ 时最大}
{$t=2\Us$ 时最小}
{$t=8.5\Us$ 时最大}
{$t=8.5\Us$ 时最小}

%% answer: AD
%% memo:
\memoanswer{
	本题原卷及材料中未提供详解，待补充。
}

\item
%% number: 7
%% paperName: 2015年江苏理综第 7 题 4 分
%% typeId: 2
%% type: 多选题
%% chapter: 电场
%% point: 带电粒子的偏转
%% method: 无
%% score: 4
%% degree: 450
%% duplicateId: 0
%% body:
一带正电的小球向右水平抛入范围足够大的匀强电场，电场方向水平向左。不计空气阻力，则小球\xzanswer{BC}

\onepicture[h=2.6cm, align=right]{figs/07.png}

\fourchoices[answer=BC]
{做直线运动}
{做曲线运动}
{速率先减小后增大}
{速率先增大后减小}

%% answer: BC
%% memo:
\memoanswer{
	本题原卷及材料中未提供详解，待补充。
}

\item
%% number: 8
%% paperName: 2015年江苏理综第 8 题 4 分
%% typeId: 2
%% type: 多选题
%% chapter: 电场
%% point: 电场线
%% method: 无
%% score: 4
%% degree: 450
%% duplicateId: 0
%% body:
两个相同的负电荷和一个正电荷附近的电场线分布如图所示。$c$ 是两负电荷连线的中点，$d$ 点在正电荷的正上方，$c$、$d$ 到正电荷的距离相等，则\xzanswer{ACD}

\onepicture[h=3.6cm, align=right]{figs/08.png}

\fourchoices[answer=ACD]
{a 点的电场强度比 b 点的大}
{a 点的电势比 b 点的高}
{c 点的电场强度比 d 点的大}
{c 点的电势比 d 点的低}

%% answer: ACD
%% memo:
\memoanswer{
	本题原卷及材料中未提供详解，待补充。
}

\item
%% number: 9
%% paperName: 2015年江苏理综第 9 题 4 分
%% typeId: 2
%% type: 多选题
%% chapter: 运动和力
%% point: 变加速直线运动
%% method: 无
%% score: 4
%% degree: 450
%% duplicateId: 0
%% body:
如图所示，轻质弹簧一端固定，另一端与一质量为 $m$、套在粗糙竖直固定杆 A 处的圆环相连，弹簧水平且处于原长。圆环从 A 处由静止开始下滑，经过 B 处的速度最大，到达 C 处的速度为零，$AC=h$。圆环在 C 处获得一竖直向上的速度 $v$，恰好能回到 A。弹簧始终在弹性限度内，重力加速度为 $g$。则圆环\xzanswer{BD}

\onepicture[h=3.6cm, align=right]{figs/09.png}

\fourchoices[answer=BD]
{下滑过程中，加速度一直减小}
{下滑过程中，克服摩擦力做的功为 $\frac{1}{4}mv^{2}$}
{在 C 处，弹簧的弹性势能为 $\frac{1}{4}mv^{2}-mgh$}
{上滑经过 B 的速度大于下滑经过 B 的速度}

%% answer: BD
%% memo:
\memoanswer{
	本题原卷及材料中未提供详解，待补充。
}

\item
%% number: 10
%% paperName: 2015年江苏理综第 10 题 8 分
%% typeId: 4
%% type: 实验题
%% chapter: 电路
%% point: 测电动势与内阻
%% method: 无
%% score: 8
%% degree: 650
%% duplicateId: 0
%% body:
小明利用如图所示的实验装置测量一干电池的电动势和内阻。

\onepicture[h=4.0cm, align=right]{figs/10a.png}

（1）图中电流表的示数为\tkanswer[1.4cm]{0.44}$\UA$；

（2）调节滑动变阻器，电压表和电流表的示数记录如下：

\begin{center}
\begin{tabular}{|c|c|c|c|c|c|}
\hline
$U$（$\UV$） & 1.45 & 1.36 & 1.27 & 1.16 & 1.06 \\
\hline
$I$（$\UA$） & 0.12 & 0.20 & 0.28 & 0.36 & 0.44 \\
\hline
\end{tabular}
\end{center}

请根据表中的数据，在答题卡的方格纸上作出 $U-I$ 图线。

\onepicture[h=5.0cm]{figs/10b.png}

由图线求得：电动势 $E=$\tkanswer[1.6cm]{1.60}$\UV$；内阻 $r=$\tkanswer[1.6cm]{1.2}$\UO$；

（3）实验时，小明进行了多次测量，花费了较长时间，测量期间一直保持电路闭合。其实，从实验误差考虑，这样的操作不妥，因为\tkanswer[5.5cm]{干电池长时间使用后，电动势和内阻会发生变化，导致实验误差增大}。

%% answer:
\jdanswer{
（1）$0.44\UA$；
（2）$E=1.60\UV$（1.58~1.62 均可），$r=1.2\UO$（1.18~1.26 均可）；
（3）干电池长时间使用后，电动势和内阻会发生变化，导致实验误差增大。
}
%% memo:
\memoanswer{
	本题原卷及材料中未提供详解，待补充。
}

\item
%% number: 11
%% paperName: 2015年江苏理综第 11 题 10 分
%% typeId: 4
%% type: 实验题
%% chapter: 电磁感应
%% point: 电磁感应中的变加速
%% method: 无
%% score: 10
%% degree: 450
%% duplicateId: 0
%% body:
某同学探究小磁铁在铜管中下落时受电磁阻尼作用的运动规律。实验装置如图所示，打点计时器的电源为 $50\UHz$ 的交流电。

\onepicture[h=4.4cm, align=right]{figs/11a.png}

（1）下列实验操作中，不正确的有\tkanswer[2.0cm]{CD}。

（A）将铜管竖直地固定在限位孔的正下方

（B）纸带穿过限位孔，压在复写纸下面

（C）用手捏紧磁铁保持静止，然后轻轻地松开让磁铁下落

（D）在磁铁下落的同时接通打点计时器的电源

（2）该同学按正确的步骤进行实验（记为“实验①”），将磁铁从管口处释放，打出一条纸带，取开始下落的一段，确定一合适的点为 O 点，每隔一个计时点取一个计数点，标为 1，2，$\ldots$，8。用刻度尺量出各计数点的相邻两计时点到 O 点的距离，记录在纸带上，如图所示。

\onepicture[h=1.6cm]{figs/11b.png}

计算相邻计时点间的平均速度 $\bar v$，粗略地表示各计数点的速度，抄入下表。请将表中的数据补充完整。

\begin{center}
\begin{tabular}{|c|c|c|c|c|c|c|c|c|}
\hline
位置 & 1 & 2 & 3 & 4 & 5 & 6 & 7 & 8 \\
\hline
$\bar v$（$\Ucms$） & 24.5 & 33.8 & 37.8 & \tkanswer[1.2cm]{39.0} & 39.5 & 39.8 & 39.8 & 39.8 \\
\hline
\end{tabular}
\end{center}

（3）分析上表的实验数据可知：在这段纸带记录的时间内，磁铁运动速度的变化情况是\tkanswer[4.0cm]{逐渐增大到 $39.8\Ucms$}；磁铁受到阻尼作用的变化情况是\tkanswer[4.0cm]{逐渐增大到等于重力}。

（4）该同学将装置中的铜管更换为相同尺寸的塑料管，重复上述实验操作（记为“实验②”），结果表明磁铁下落的运动规律与自由落体运动规律几乎相同。请问实验②是为了说明什么？对比实验①和②的结果可得到什么结论？

\tkanswer[10cm]{为了说明磁铁在塑料管中几乎不受阻尼作用。磁铁在铜管中受到的阻尼作用主要是电磁阻尼作用。}

%% answer:
\jdanswer{
（1）CD；
（2）39.0；
（3）逐渐增大到 $39.8\Ucms$，逐渐增大到等于重力；
（4）为了说明磁铁在塑料管中几乎不受阻尼作用。磁铁在铜管中受到的阻尼作用主要是电磁阻尼作用。
}
%% memo:
\memoanswer{
	本题原卷及材料中未提供详解，待补充。
}

\item
%% number: 12
%% paperName: 2015年江苏理综第 12 题 4 分
%% typeId: 6
%% type: 计算题
%% chapter: 原子和原子核
%% point: 结合能
%% method: 无
%% score: 4
%% degree: 650
%% duplicateId: 0
%% body:
【选做题】本题包括 A、B、C 三小题，请选定其中两小题，并在相应的答题区域内作答。若多做，则按 A、B 两小题评分。

\textbf{A．}（选修 3-3）（12 分）

\begin{enumerate}
	\item
	对下列几种固体物质的认识，正确的有\xzanswer{CD}

	\fourchoices[answer=CD]
	{食盐熔化过程中，温度保持不变，说明食盐是晶体}
	{烧热的针尖接触涂有蜂蜡薄层的云母片背面，熔化的蜂蜡呈椭圆形，说明蜂蜡是晶体}
	{天然石英表现为各向异性，是由于该物质的微粒在空间的排列不规则}
	{石墨和金刚石的物理性质不同，是由于组成它们的物质微粒排列结构不同}

	\item
	在装有食品的包装袋中充入氮气，可以起到保质作用。某厂家为检测包装袋的密封性，在包装袋中充满一定量的氮气，然后密封进行加压测试。测试时，对包装袋缓慢地施加压力。将袋内的氮气视为理想气体，则加压测试过程中，包装袋内壁单位面积上所受气体分子撞击的作用力\tkanswer[1.6cm]{增大}（选填“增大”、“减小”或“不变”），包装袋内氮气的内能\tkanswer[1.6cm]{不变}（选填“增大”、“减小”或“不变”）。

	\item
	给某包装袋充入氮气后密封，在室温下，袋中气体压强为 1 个标准大气压、体积为 $1\UL$。将其缓慢压缩到压强为 2 个标准大气压时，气体的体积变为 $0.45\UL$。请通过计算判断该包装袋是否漏气。

\end{enumerate}

\textbf{B．}（选修 3-4）（12 分）

\begin{enumerate}
	\item
	一渔船向鱼群发出超声波，若鱼群正向渔船靠近，则被鱼群反射回来的超声波与发出的超声波相比\xzanswer{BC}

	\fourchoices[answer=BC]
	{波速变大}
	{波速不变}
	{频率变高}
	{频率不变}

	\item
	用 $2\times10^{6}\UHz$ 的超声波检查胆结石，该超声波在结石和胆汁中的波速分别为 $2250\Ums$ 和 $1500\Ums$，则该超声波在结石中的波长是胆汁中的\tkanswer[1.2cm]{1.5}倍。用超声波检查胆结石是因为超声波的波长较短，遇到结石时\tkanswer[1.6cm]{不容易}（选填“容易”或“不容易”）发生衍射。

	\item
	人造树脂是常用的眼镜镜片材料。如图所示，光线射在一人造树脂立方体上，经折射后，射在桌面上的 P 点。已知光线的入射角为 $30^{\circ}$，$OA=5\Ucm$，$AB=20\Ucm$，$BP=12\Ucm$，求该人造树脂材料的折射率 $n$。

	\onepicture[h=3.4cm, align=right]{figs/12.png}

\end{enumerate}

\textbf{C．}（选修 3-5）（12 分）

\begin{enumerate}
	\item
	波粒二象性是微观世界的基本特征，以下说法正确的有\xzanswer{AB}

	\fourchoices[answer=AB]
	{光电效应现象揭示了光的粒子性}
	{热中子束射到晶体上产生衍射图样说明中子具有波动性}
	{黑体辐射的实验规律可用光的波动性解释}
	{动能相等的质子和电子，它们的德布罗意波长也相等}

	\item
	核电站利用原子核链式反应放出的巨大能量进行发电，$\ce{^{235}_{92}U}$ 是核电站常用的核燃料。$\ce{^{235}_{92}U}$ 受一个中子轰击后裂变成 $\ce{^{144}_{56}Ba}$ 和 $\ce{^{89}_{36}Kr}$ 两部分，并产生\tkanswer[1.2cm]{3}个中子。要使链式反应发生，裂变物质的体积要\tkanswer[1.6cm]{大于}（选填“大于”或“小于”）它的临界体积。

	\item
	取质子的质量 $m_{\mathrm{p}}=1.6726\times10^{-27}\Ukg$，中子的质量 $m_{\mathrm{n}}=1.6749\times10^{-27}\Ukg$，$\alpha$ 粒子的质量 $m_{\alpha}=6.6467\times10^{-27}\Ukg$，光速 $c=3.0\times10^{8}\Ums$。请计算 $\alpha$ 粒子的结合能。（计算结果保留两位有效数字）

\end{enumerate}

%% answer:
\jdanswer{
\begin{enumerate}
	\item
	A：（1）CD；（2）增大，不变；（3）因为 $0.45\UL<0.5\UL$，故包装袋漏气。

	\item
	B：（1）BC；（2）1.5，不容易；（3）$n=\frac{\sqrt{449}}{14}$（或 $n\approx1.5$）。

	\item
	C：（1）AB；（2）3，大于；（3）$\Delta E=4.3\times10^{-12}\UJ$。
\end{enumerate}
}
%% memo:
\memoanswer{
A．（1）（2）本题原卷及材料中未提供详解，待补充。

（3）若不漏气，设加压后的体积为 $V_{1}$，由等温过程得 $p_{0}V_{0}=p_{1}V_{1}$，代入数据得 $V_{1}=0.5\UL$。因为 $0.45\UL<0.5\UL$，故包装袋漏气。

B．（1）（2）本题原卷及材料中未提供详解，待补充。

（3）设折射角为 $\gamma$，由折射定律 $\sin30^{\circ}=n\sin\gamma$；由几何关系知 $\sin\gamma=\frac{PB-OA}{OP}$，且 $OP=\sqrt{(PB-OA)^{2}+AB^{2}}$。代入数据解得 $n=\frac{\sqrt{449}}{14}$（或 $n\approx1.5$）。

C．（1）（2）本题原卷及材料中未提供详解，待补充。

（3）组成 $\alpha$ 粒子的核子与 $\alpha$ 粒子的质量差 $\Delta m=(2m_{\mathrm{p}}+2m_{\mathrm{n}})-m_{\alpha}$，结合能 $\Delta E=\Delta mc^{2}$，代入数据得 $\Delta E=4.3\times10^{-12}\UJ$。
}

\item
%% number: 13
%% paperName: 2015年江苏理综第 13 题 15 分
%% typeId: 6
%% type: 计算题
%% chapter: 电磁感应
%% point: 法拉第电磁感应定律
%% method: 无
%% score: 15
%% degree: 650
%% duplicateId: 0
%% body:
做磁共振（MRI）检查时，对人体施加的磁场发生变化时会在肌肉组织中产生感应电流。某同学为了估算该感应电流对肌肉组织的影响，将包裹在骨骼上的一圈肌肉组织等效成单匝线圈，线圈的半径 $r=5.0\Ucm$，线圈导线的截面积 $A=0.80\Ucmq$，电阻率 $\rho=1.5\UOm$。如图所示，匀强磁场方向与线圈平面垂直，若磁感应强度 $B$ 在 $0.3\Us$ 内从 $1.5\UT$ 均匀地减为零，求：（计算结果保留一位有效数字）

\onepicture[h=2.2cm, align=right]{figs/13.png}

\begin{enumerate}
	\item
	该圈肌肉组织的电阻 $R$；

	\item
	该圈肌肉组织中的感应电动势 $E$；

	\item
	$0.3\Us$ 内该圈肌肉组织中产生的热量 $Q$。
\end{enumerate}

%% answer:
\jdanswer{
\begin{enumerate}
	\item
	$R=6\times10^{3}\UO$

	\item
	$E=4\times10^{-2}\UV$

	\item
	$Q=8\times10^{-8}\UJ$
\end{enumerate}
}
%% memo:
\memoanswer{
（1）由电阻定律得 $R=\rho\frac{2\pi r}{A}$，代入数据得 $R=6\times10^{3}\UO$。

（2）感应电动势 $E=\frac{\Delta B\pi r^{2}}{\Delta t}$，代入数据得 $E=4\times10^{-2}\UV$。

（3）由焦耳定律得 $Q=\frac{E^{2}}{R}\Delta t$，代入数据得 $Q=8\times10^{-8}\UJ$。
}

\item
%% number: 14
%% paperName: 2015年江苏理综第 14 题 16 分
%% typeId: 6
%% type: 计算题
%% chapter: 曲线运动
%% point: 圆周运动的应用
%% method: 无
%% score: 16
%% degree: 450
%% duplicateId: 0
%% body:
一转动装置如图所示，四根轻杆 OA、OC、AB 和 CB 与两小球及一小环通过铰链连接，轻杆长均为 $l$，球和环的质量均为 $m$，O 端固定在竖直的轻质转轴上。套在转轴上的轻质弹簧连接在 O 与小环之间。原长为 $L$。装置静止时，弹簧长为 $3L/2$。转动该装置并缓慢增大转速，小环缓慢上升。弹簧始终在弹性限度内，忽略一切摩擦和空气阻力，重力加速度为 $g$。求

\onepicture[h=3.6cm, align=right]{figs/14.png}

\begin{enumerate}
	\item
	弹簧的劲度系数 $k$；

	\item
	AB 杆中弹力为零时，装置转动的角速度 $\omega_{0}$；

	\item
	弹簧长度从 $3L/2$ 缓慢缩短为 $L/2$ 的过程中，外界对转动装置所做的功 $W$。
\end{enumerate}

%% answer:
\jdanswer{
\begin{enumerate}
	\item
	$k=\frac{4mg}{L}$

	\item
	$\omega_{0}=\sqrt{\frac{8g}{5L}}$

	\item
	$W=mgL+\frac{16mgl^{2}}{L}$
\end{enumerate}
}
%% memo:
\memoanswer{
（1）装置静止时，设 OA、AB 杆中的弹力分别为 $F_{1}$、$T_{1}$，OA 杆与转轴的夹角为 $\theta_{1}$。

小环受到弹簧的弹力 $F_{\text{弹}1}=kL/2$，

小环受力平衡 $F_{\text{弹}1}=mg+2T_{1}\cos\theta_{1}$，

小球受力平衡 $F_{1}\cos\theta_{1}+T_{1}\cos\theta_{1}=mg$；$F_{1}\sin\theta_{1}=T_{1}\sin\theta_{1}$，

解得 $k=\frac{4mg}{L}$。

（2）设 OA、AB 杆中的弹力分别为 $F_{2}$、$T_{2}$，OA 杆与转轴的夹角为 $\theta_{2}$，弹簧长度为 $x$。小环受到弹簧的弹力 $F_{\text{弹}2}=k(x-L)$，

小环受力平衡 $F_{\text{弹}2}=mg$ 得 $x=5L/4$，

对小球 $F_{2}\cos\theta_{2}=mg$；$F_{2}\sin\theta_{2}=m\omega_{0}^{2}l\sin\theta_{2}$ 且 $\cos\theta_{2}=x/2L$，

解得 $\omega_{0}=\sqrt{\frac{8g}{5L}}$。

（3）弹簧长度为 $L/2$ 时，设 OA、AB 杆中的弹力分别为 $F_{3}$、$T_{3}$，OA 杆与弹簧的夹角为 $\theta_{3}$。

小环受到弹簧的弹力 $F_{\text{弹}3}=kL/2$，

小环受力平衡 $2T_{3}\cos\theta_{3}=mg+F_{\text{弹}3}$ 且 $\cos\theta_{3}=L/4l$，

对小球 $F_{3}\cos\theta_{3}=T_{3}\cos\theta_{3}+mg$；$F_{3}\sin\theta_{3}+T_{3}\sin\theta_{3}=m\omega_{3}^{2}l\sin\theta_{3}$，

解得 $\omega_{3}=\sqrt{\frac{16g}{L}}$。

整个过程弹簧弹性势能变化为零，则弹力做的功为零，由动能定理

$W-mg\left(\frac{3L}{2}-\frac{L}{2}\right)-2mg\left(\frac{3L}{4}-\frac{L}{4}\right)=2\times\frac{1}{2}m(\omega_{3}l\sin\theta_{3})^{2}$，

解得 $W=mgL+\frac{16mgl^{2}}{L}$。
}

\item
%% number: 15
%% paperName: 2015年江苏理综第 15 题 16 分
%% typeId: 6
%% type: 计算题
%% chapter: 磁场
%% point: 带电粒子在磁场中的运动
%% method: 无
%% score: 16
%% degree: 250
%% duplicateId: 0
%% body:
一台质谱仪的工作原理如图所示，电荷量均为 $+q$、质量不同的离子飘入电压为 $U_{0}$ 的加速电场，其初速度几乎为零。这些离子经加速后通过狭缝 O 沿着与磁场垂直的方向进入磁感应强度为 $B$ 的匀强磁场，最后打在底片上。已知放置底片的区域 $MN=L$，且 $OM=L$。某次测量发现 MN 中左侧 2/3 区域 MQ 损坏，检测不到离子，但右侧 1/3 区域 QN 仍能正常检测到离子。在适当调节加速电压后，原本打在 MQ 的离子即可在 QN 检测到。

\onepicture[h=3.4cm, align=right]{figs/15.png}

\begin{enumerate}
	\item
	求原本打在 MN 中点 P 的离子质量 $m$；

	\item
	为使原本打在 P 的离子能打在 QN 区域，求加速电压 $U$ 的调节范围；

	\item
	为了在 QN 区域将原本打在 MQ 区域的所有离子检测完整，求需要调节 $U$ 的最少次数。（取 $\lg2=0.301$，$\lg3=0.477$，$\lg5=0.699$）
\end{enumerate}

%% answer:
\jdanswer{
\begin{enumerate}
	\item
	$m=\frac{9qB^{2}L^{2}}{32U_{0}}$

	\item
	$\frac{100}{81}U_{0}\le U\le\frac{16}{9}U_{0}$

	\item
	最少次数为 3 次
\end{enumerate}
}
%% memo:
\memoanswer{
（1）离子在电场中加速 $qU_{0}=\frac{1}{2}mv^{2}$，

在磁场中做匀速圆周运动 $qvb=m\frac{v^{2}}{r}$，

解得 $r=\frac{1}{B}\sqrt{\frac{2mU_{0}}{q}}$，

代入 $r_{0}=\frac{3}{4}L$，解得 $m=\frac{9qB^{2}L^{2}}{32U_{0}}$。

（2）由（1）知，$U=\frac{16U_{0}r^{2}}{9L^{2}}$。离子打在 Q 点 $r=\frac{5}{6}L$，$U=\frac{100U_{0}}{81}$；

离子打在 N 点 $r=L$，$U=\frac{16U_{0}}{9}$。则电压的范围 $\frac{100}{81}U_{0}\le U\le\frac{16}{9}U_{0}$。

（3）由（1）可知，$r\propto\sqrt{U}$。

由题意知，第 1 次调节电压到 $U_{1}$，使原本 Q 点的离子打在 N 点 $\frac{L}{\frac{5}{6}L}=\frac{\sqrt{U_{1}}}{\sqrt{U_{0}}}$，

此时，原本半径为 $r_{1}$ 的打在 $Q_{1}$ 的离子打在 Q 上 $\frac{\frac{5}{6}L}{r_{1}}=\frac{\sqrt{U_{1}}}{\sqrt{U_{0}}}$，

解得 $r_{1}=\left(\frac{5}{6}\right)^{2}L$。

第 2 次调节电压到 $U_{2}$，原本打在 $Q_{1}$ 的离子打在 N 点，原本半径为 $r_{2}$ 的打在 $Q_{2}$ 的离子打在 Q 上，则

$\frac{L}{r_{1}}=\frac{\sqrt{U_{2}}}{\sqrt{U_{0}}}$，$\frac{\frac{5}{6}L}{r_{2}}=\frac{\sqrt{U_{2}}}{\sqrt{U_{0}}}$，解得 $r_{2}=\left(\frac{5}{6}\right)^{3}L$。

同理，第 $n$ 次调节电压，有 $r_{n}=\left(\frac{5}{6}\right)^{n+1}L$。

检测完整，有 $r_{n}\le\frac{L}{2}$，解得 $n\ge\frac{\lg2}{\lg\left(\frac{6}{5}\right)}-1\approx2.8$。

最少次数为 3 次。
}

\end{enumerate}

\end{document}
